\begin{proposition}[A finite-field rank conjecture controls algebraic homological projectors]
\label{c:chow-unramified-rank-reduction}
\status{Proved conditional implication; no converse is asserted}
Fix a prime $p$. Assume independence of $\ell$ of homological
equivalence for rational cycles on every smooth projective variety
over $\overline{\F}_p$. Equivalently, assume independence of the
ranks of all rational algebraic correspondences on their individual
cohomology groups; this is the exact equivalence of
\citet[Theorem~1, Definition~4.1 and Theorem~4.6]{c_vandobben2020}.

Let $K$ be a number field, $v\mid p$, and let $X/K$ be smooth
projective of pure dimension $d$. Let
$e\in\CH^d(X\times_KX)_\Q$ be a rational algebraic correspondence
whose Betti action is idempotent and vanishes outside degree $i$.
For an integer $m$, put
$M_\ell=\operatorname{im}(e:H^i(X_{\bar K},\Q_\ell)(m)
\to H^i(X_{\bar K},\Q_\ell)(m))$.
Betti--\'etale comparison gives the same homogeneous idempotent
condition for every $\ell$. We require no identity $e^2=e$ in
the Chow ring; the assertion therefore covers algebraic
homological projectors as well as rational Chow motives.
Then the ranks of all powers of the local monodromy operator on
$M_\ell$ are independent of $\ell\ne p$.
Moreover,
\[
 M_{\ell_0}\text{ unramified at }v\text{ for one }\ell_0\ne p
 \quad\Longrightarrow\quad
 M_\ell\text{ unramified at }v\text{ for every }\ell\ne p.
\]
This implication does not assume weight--monodromy.
\end{proposition}

\begin{proof}
We work over the local field $K_v$ and first suppress the Tate twist:
for $\ell\ne p$ it has trivial inertia and zero monodromy.
There are three linear-algebra and spectral-sequence points to verify.

\emph{1. Ranks on cohomology are finite block-matrix ranks.}
Let $f:C^\bullet\to D^\bullet$ be a map of finite complexes over a
field. For every $a$,
\begin{equation}
\label{eq:c-induced-cohomology-rank}
 \operatorname{rank}H^a(f)
 =\operatorname{rank}
 \begin{pmatrix}
  d_C^a&0\\ f^a&d_D^{a-1}
 \end{pmatrix}
 -\operatorname{rank}d_C^a-\operatorname{rank}d_D^{a-1}.
\end{equation}
Indeed, projection of the block image onto $C^{a+1}$ has image
$\operatorname{im}d_C^a$ and kernel
$f^a(\ker d_C^a)+\operatorname{im}d_D^{a-1}$.
Dividing that kernel by $\operatorname{im}d_D^{a-1}$ gives exactly
the image on cohomology, since $f$ carries boundaries to boundaries.

The rank conjecture applies to these block matrices whenever their
entries are algebraic correspondences between finite sums of twisted
stratum cohomology groups of one weight. To check this assertion
explicitly, write a summand as $H^{j_s}(Z_s)(t_s)$ with
$j_s-2t_s=w$. Choose $T\geq\max_s t_s$ and put
$r_s=T-t_s\geq0$. The summand identifies with the part of
\[
 H^{w+2T}(Z_s\times\mathbf P^{r_s})(T)
\]
selected by the algebraic top-degree projector of the projective
space. Exterior products of the original correspondences with the
projective-space correspondences implement every block.
Choose a common dimension $D$ at least as large as every
$\dim Z_s+r_s$, and further multiply the $s$th component by
$\mathbf P^{D-\dim Z_s-r_s}$ with its algebraic degree-zero
projector. This leaves the selected cohomological degree unchanged
and makes all components equidimensional of dimension $D$.
Disjoint unions implement finite direct sums, and the top-degree
and degree-zero projectors kill the unwanted projective-space
components. Thus the block map is a single codimension-$D$
rational algebraic correspondence
on a fixed cohomological degree of smooth projective schemes.
Its rank is independent of $\ell$ by the stated conjecture.
No algebraic projector onto a kernel of a differential is required.

\emph{2. A semistable projector has common algebraic $E_1$ maps.}
Choose a finite Galois extension $L/K_v$, a projective strictly
semistable scheme $\mathcal Y/\mathcal O_L$, and a proper
generically finite map $f:Y\to X_L$ of positive degree $\delta$.
Such an alteration is supplied by
\citet[Corollary~1.2.6, author p.~16]{c_saito2003weights}.
For geometric components, take disjoint unions of alterations and,
if necessary, repeat components to make the generic degree common.
The correspondence
\[
                 e_Y=\delta^{-1}f^* e f_*
\]
is algebraic and acts idempotently on cohomology. The identities $f_*f^*=\delta$ and
the projection formula identify its realized image with that of
$e$, compatibly with Galois action and monodromy.
Passing from $K_v$ to $L$ multiplies the normalized monodromy
operator by the nonzero ramification index, so none of its ranks
changes.

Use Saito's convention that $Y^{(s)}$ is the disjoint union of the
$(s+1)$-fold intersections of geometric components. The weight
spectral sequence is
\[
 E_{1,\ell}^{a,b}
 =\bigoplus_{j\geq\max(0,-a)}
 H^{b-2j}(Y^{(a+2j)}_{\overline k},\Q_\ell)(-j)
 \ \Longrightarrow\ H^{a+b}(Y_{\overline L},\Q_\ell).
\]
Its differentials are fixed signed restriction and Gysin
correspondences. The action of $e_Y$ on this spectral sequence
has $E_1$ components induced by fixed rational algebraic
correspondences on the strata, independent of $\ell$.
This is precisely
\citet[Lemma~2.4.2, Corollary~2.4.3 and Proposition~2.4.4,
complete proofs, author pp.~34--36]{c_saito2003weights}:
one resolves the closure of a cycle on the regular semistable
product blowup by vector bundles, restricts their Chern characters
to the strata, and divides by the fixed factorial.
The resulting rational cycles are chosen before $\ell$.
The $E_1$ endomorphism need not be idempotent; its induced action
on $E_2$ and the abutment is the required projector.

On $E_1$, the monodromy map
\[
 N:E_{1,\ell}^{a,b}\longrightarrow
       E_{1,\ell}^{a+2,b-2}(-1)
\]
is the identity on matching summands and zero on the others.
It sends $(a,j)$ to $(a+2,j-1)$ and hence preserves the stratum
index $a+2j$. Consequently the stratum correspondence blocks
commute with this map. The assertion for $N$ follows from
\citet[Corollary~2.2.4(2), author pp.~21--22]{c_saito2003weights}:
the displayed map for $T-1$ differs on the graded pieces from
$N$ only by its nonzero tame scalar. Higher logarithm terms drop
the filtration further and do not change this leading map.

Fix $b$ and an integer $r\geq1$. Regard
$C^a=E_{1,\ell}^{a,b}$ and
$D^a=E_{1,\ell}^{a+2r,b-2r}(-r)$ as complexes in $a$.
All their summands have weight $b$, and
$N^r e_Y:C^\bullet\to D^\bullet$ has rational algebraic
stratum blocks. Formula~\eqref{eq:c-induced-cohomology-rank}
therefore shows that the rank of its induced $E_2$ map is
independent of $\ell$.
The spectral sequence degenerates at $E_2$, since its $E_1$
rows are pure of their distinct weights. Thus these are the ranks
of all graded maps
\[
 N^r e_Y:\operatorname{Gr}^W_b H^i(Y)(r)
       \longrightarrow\operatorname{Gr}^W_{b-2r}H^i(Y),
\]
or, equivalently, of their simultaneous untwisted form with
$(-r)$ on the target.

\emph{3. Graded monodromy detects actual monodromy here.}
Let $F$ be a geometric Frobenius lift over $L$ and let $q_L$
be its residue cardinality. Each $\operatorname{Gr}^W_b H^i(Y)$
is pure of weight $b$, by
\citet[Lemma~3.7, author p.~42]{c_saito2003weights}.
Distinct weights have disjoint Frobenius spectra. Consequently
the filtration has a unique $F$-equivariant splitting by the
generalized Frobenius weight spaces
\[
             H^i(Y)=\bigoplus_b V_b,\qquad
             W_b H^i(Y)=\bigoplus_{c\leq b}V_c.
\]
This follows from primary decomposition for $F$ and does not
assume that $F$ is semisimple. The relation
$FNF^{-1}=q_L^{-1}N$ carries $V_b$ into $V_{b-2}$,
while $e_Y$ preserves every $V_b$. Therefore
\[
 \operatorname{rank}(N^r e_Y\mid H^i(Y))
 =\sum_b\operatorname{rank}
       (\operatorname{Gr}^W_b(N^r e_Y)).
\]
In particular zero on all graded pieces is equivalent to zero
on the entire space. Because $e_Y$ is a projector commuting with
$N$, the left side is the rank of $N^r$ on its image.
This proves the first assertion, including the independence
of vanishing of $N$.

Finally suppose $M_{\ell_0}$ is unramified. Its monodromy
vanishes, hence the preceding argument makes monodromy vanish
on $M_\ell$ for every $\ell\ne p$. Strict semistability of
$\mathcal Y$ then makes $I_L$ act trivially on every such image.
Thus the original inertia actions factor through one finite
quotient of $I_{K_v}$.
By Saito's Theorem~0.1, the alternating trace of $e\sigma$ is
rational and independent of $\ell$ for every
$\sigma\in I_{K_v}$. Homogeneity of $e$ identifies this with
$(-1)^i\operatorname{Tr}(\sigma\mid M_\ell)$; the Tate twist
is trivial on inertia. At $\ell_0$ the trace is the common
dimension of $M$, so it is that dimension for every $\ell$.
Averaging these traces over the finite inertia quotient gives
the full dimension for its invariant subspace. Finite-group
representations in characteristic zero are semisimple, hence
the inertia action is trivial for every $\ell$.
\end{proof}

\begin{remark}[The precise conditional and categorical boundary]
\label{c:chow-rank-reduction-scope}
\status{Proved sufficient-conjecture reduction; the general equivalence is open}
The hypothesis is a published, independently formulated
independence-of-homological-equivalence conjecture over
$\overline{\F}_p$, not a renamed version of C.6.
The result proves only its implication to the stated C.6 assertion.
It does not prove the converse: realizing an arbitrary finite-field
correspondence as local monodromy of a number-field variety would
require an additional geometric construction which is not supplied.
The representative must be algebraic and homogeneous, but only
its cohomological action must be idempotent. Inspecting the proof,
idempotence is used on the realization and its $E_2$ graded pieces,
never in the Chow group. No extension to an arbitrary
absolute-Hodge, motivated, or numerical projector is inferred
from Saito's algebraic-correspondence construction alone.
The following theorem treats motivated projectors by an additional
finite algebraic-and-Lefschetz construction.
Thus this proposition gives a second exact sufficient hypothesis,
alongside the purity reduction, but does not satisfy the requested
general-motive proof or a two-way named-conjecture equivalence.
\end{remark}
